\chapter*{Abstract}
\chaptermark{Abstract}
\addcontentsline{toc}{chapter}{Abstract}  

%This chapter should contain three things.

%\begin{itemize}
    %\item A copy of all the information on the title page of your master's thesis. This includes things like the name of your master's thesis and your advisors.
    %\item A one-paragraph description of your master's thesis. This should be 15 to 20 lines long. This should include the context of your master's thesis, the problem statement of your master's thesis. The results of your master's thesis, and the evaluation of the work.
    %\item Five keywords that describe the subject best.
%\end{itemize}

%The chapter should be one page at most.

\textbf{Title:} A Congestion-Aware Two-Sided Job Matching Framework\\
\textbf{Author:} Pharrell Tratsaert\\
\textbf{Student number:} 02111011\\
\textbf{Supervisor:} Prof. dr. Bo Kang\\
\textbf{Counsellor:} Prof. dr. Tijl De Bie\\
\textbf{Programme:} Master of Science in Computer Science Engineering\\
\textbf{Academic year:} 2025-2026
\bigskip
\bigskip

The job market is a two-sided platform where jobseekers and recruiters
look for the perfect match. A structural problem in
such markets is congestion: when a recommender system recommends the same
popular jobs to many users, applications concentrate on a small subset
of postings, leaving the majority of jobs starved while a
small number is overwhelmed, and causing most candidates to be
overlooked by recruiters entirely because of the large number of applications for the popular jobs. This dissertation proposes a
congestion-aware two-phase matching framework to reduce this congestion while
preserving fair and efficient matches. On the user side, a Two-Tower
neural architecture generates personalised job recommendation lists. On
the recruiter side, a lightweight XGBoost model scores candidate
suitability, trained on synthetic labels generated by a large language
model using Chain-of-Thought prompting. Both components are integrated
into a hybrid congestion-balancing strategy. For the job recommendation stage, a combination of bilateral scoring and market-maker penalties is used to adjust recommendation scores based on the current congestion level of each job and to improve suitability of a recommendation, while for the shortlisting stage, a candidate-level penalty is applied to down-weight candidates who have already been shortlisted by many recruiters. This creates a dynamic feedback loop that steers attention towards less popular jobs and candidates, while still prioritising relevance. The framework is
evaluated through counterfactual simulations on real application logs
from the CareerBuilder dataset, measuring applications-per-job and
shortlists-per-candidate distributions, Gini coefficients, and
model-based match probabilities across a held-out deployment window.
The results show that the full two-phase pipeline substantially reduces congestion on both sides of the market compared to an uncorrected Two-Tower recommender. On the job side, the Gini coefficient drops from 0.9625 to 0.7130 and the fraction of starved jobs decreases by 27 percentage points. On the candidate side, the Gini drops from 0.7896 to 0.4967 and the fraction of starved candidates falls by 76 percentage points. In terms of match quality, the congestion reduction comes at a modest cost to user-side relevance: the average Two-Tower probability decreases by 9 percentage points. However, both recruiter-side signals improve substantially: the average recruiter suitability over all applied pairs increases by 31 percentage points, and the average suitability of shortlisted candidates increases by 19 percentage points. This demonstrates that while the framework steers users toward less popular jobs, it simultaneously improves the quality of matches from the recruiter's perspective, benefiting both sides of the market.

\textbf{Keywords:} recommender systems; job matching; two-sided markets; market congestion; machine learning
